% !TeX root = ../papers/nc-p3.tex
\subsection{特異三次曲線の三生成基底と楕円枝への平滑化}
\label{sec:singular-tangent-smoothing}
引き続き四辺が有限で三次元接点を持つ枝を扱う。
三生成基底の点対応が既約特異三次曲線のグラフの場合も決定できる。

\begin{theorem}\label{thm:nodal-tangent-windows}
三生成基底の対応が射影平面の射影自己同型のグラフ、または尖点三次曲線の自己同型の
グラフなら、この有限四辺の枝に入らない。既約節点三次曲線のグラフであれば、
正規化の座標を $p(t)=[t:t^2:1+t^3]$ と選んで移動は $t\mapsto\lambda t$ であり、
$\lambda=-1$ または $\lambda^2=-1$ である。基底の三関係は
\[
 xy-\lambda yx,\qquad yz-\lambda zy+\beta x^2,\qquad
 zx-\lambda xz+\beta y^2,\qquad \beta=(\lambda^3-1)/\lambda .
\]
どの基底にも九次元の全補正空間 $P_B$ がある。
有限条件は $v=(0,0,1)$ で $(d_0(x)(v,v),d_0(y)(v,v))\ne(0,0)$ であり、
独立な二本の線型式の共通核を除く。指定点は一意である。
各反転対 $\{\lambda,\lambda^{-1}\}$ の有限な射影補正は、
$t\mapsto ct$ と $t\mapsto c/t$（$c^3=1$）のうち移動を対応させるもので割ると
添字保持同型類をちょうど一度ずつ与える。二型の同型類の族は各八次元で、
一般最小周期は順に2と4である。
\end{theorem}
\begin{proof}
平面グラフ $\sigma=[G]$ では $S$ の評価は外積 $u\wedge G^{-1}v$ で与えられる。
実際、このグラフで消える双一次式の空間は三次元である。
$p,Gp,G^2p$ で前節の二つの縮約平面は
$Gp\wedge H^*$ と $p\wedge H^*$ となる。
$Gp\not\in kp$ なら異なるため、補正の評価は $\langle p,Gp\rangle$ に入り、
そのような全ての点で第四座標を解ける。恒等の場合の固定点による除外は
定理\ref{thm:quantum-tangent-finite}の証明による。

既約三次曲線で移動が恒等なら、対角で消える双一次式は交代式だけなので
やはり平面グラフになる。一般点で非恒等なら前節の四移動の必要条件の証明は
そのまま使え、$\sigma^4$ は平面に延び、滑らかな固定点を持つことは許されない
（ここで固定点は $\sigma$ の固定点をいう）。
尖点三次曲線の正規化で滑らかな部分はアフィン直線なので、固定点のない自己同型は
非零の平行移動である。四乗も非零の平行移動で、平面三次曲線の唯一の滑らかな
変曲点を保たない。変曲点の一意性は尖点三次曲線 $y^2z=x^3$ の接線との交叉でも従う。

節点三次曲線の正規化では節点上の二点を $0,\infty$ と置く。
三次の斉次式の両端係数が所定の非零比を持つという一条件が降下条件なので、
$t$ の倍率と平面の基底を選べば $p(t)$ の表示を得る。
自己同型 $t\mapsto c/t$ は滑らかな固定点を持つので、移動は $t\mapsto\lambda t$ に限る。
$\lambda^3=1$ なら平面に延び、三次曲線には二次の方程式がないので平面グラフとなる。
$\lambda^3\ne1$ とする。$(p(t),p(\lambda t))$ を代入すると表示した三関係は零となる。
九つの評価積は $t,t^2,t^3,t^4,t^5,1+\lambda^3t^6$ を張るので核はこれで尽きる。
三次交叉は三つの関係を対応する文字で左から掛けて足した巡回テンソルであり、
前節の $\ell=r$ が成立する。従って $\lambda^4=1$ が必要となる。

逆に、この基底は量子平面 $xy=\lambda yx$ に $z$ を加えたOre拡大である。
自己同型は $(x,y)\mapsto(\lambda x,\lambda^{-1}y)$、補正は
$x\mapsto-\beta y^2$, $y\mapsto(\beta/\lambda)x^2$ であり、
$xy-\lambda yx$ に捩れた積の規則を適用すると零になる。
従って順序付き基底を持ち、$z$ の個数による随伴次数は三生成量子多項式環である。
定理\ref{thm:maximal-tangent-ore}のbar濾過の証明により標準三次元Koszul左右AS性を得る。
点対応の行列式は上の節点三次曲線を定め、節点でも行列は階数2なので
両方向の核直線によりグラフが得られる。

九次元と節点での二方向の全射性は下の六基底の同時計算で示す。
補正の節点での評価を $w_n=d_n(v,v)$ と書く。その $kv$ を法とする類は補生成元変更に依存しない。
$v$ が固定点なので、類が零なら節点の上に一次元の解がある。非零なら節点の上に解はない。
節点における式\eqref{eq:ternary-shift-derivation}の両縮約行列の核は $kv$ なので、
この非零性は次の位置と同値である。滑らかな部分では四座標の残りの条件は
\[
                 f_n(t)=\det[p(t),p(\lambda t),d_n(p(t),p(\lambda t))]=0
\]
という一変数の式である。節点の $y$ 成分が非零なら $t=0$ の近傍で
$f_n$ の一次係数が非零である。$x$ 成分が非零なら $t=\infty$ の自明化で
同様である。これは二枝の接線がそれぞれ $x,y$ 方向で、移動の微分が
$\lambda,\lambda^{-1}\ne1$ であることによる。従ってこの式は恒等零にならない。
各根の局所環は重複度を長さとする曲線状環、指定点の長さは
$20-(\deg f_n-\operatorname{ord}_0 f_n)$ であり、三次元接点は一意である。

節点を保つ正規化の変換のうち、埋込みの三次切断空間を保つのは表示した六つだけである。
反転は $\lambda$ と $\lambda^{-1}$ を交換する。定理\ref{thm:ternary-ore-linear-space}
から補正にはその変換の誘導作用を使えばよく、共通スカラーは射影化で除かれる。
従って零係数や重根でも有限群による一致が同型に必要十分である。
特殊値の最小周期は同じ射影軌道を与える $F_B^j$ の最小の正べきである。
有限な射影開集合は八次元である。
原始四乗根の型では比較に残る群の位数は3なので、固有値が下表の $F_B$ の
真のべき $j=1,2$ と射影変換として一致する元はない。
$-1$ 型で $F_B$ は全群と可換である。三乗根倍は節点評価の二方向に異なる
重みで作用するので、その位数は射影作用でも3である。
従って二面体群の射影像は忠実で中心が自明であり、非スカラーの $F_B$ に等しい元はない。
各群元との一致は閉条件なので一般最小周期が従う。
\end{proof}

\begin{theorem}\label{thm:ternary-torsion-smoothing}
上の節点基底の全ての有限窓と、定理\ref{thm:quantum-tangent-finite}の全ての有限窓は、
それぞれ同じ位数の滑らかな楕円基底を持つ窓の平坦特殊化である。
さらに二次曲線と横断的な直線の和に対応する下記の一係数基底でも同じことが成立する。
\end{theorem}
\begin{proof}
三関係を $x_{j+1}x_{j+2}+b x_{j+2}x_{j+1}+c_jx_j^2$ と書く。
$b=-\lambda$ を特別値とし、巡回順に $c_0,\ldots,c_{h-1}$ は非零、
残りは零とする（$h=0,1,2$）。対角変数変更で非零係数は1にできる。
$h=0$ は三角形、$h=2$ は前記の節点基底であり、$h=1$ では行列式が
$x_0((1-\lambda^3)x_1x_2+\lambda c_0x_0^2)$ の定数倍となる。
全基底の正則性は節点の場合と同じ量子平面上のOre表示から従う
（補正の一方または両方を零にしても関係上の導分条件は保たれる）。

実際の三次商を使う行列計算は
\path{computations/verify_singular_tangent_smoothing.py}にある。
平方係数を $1$ から $h$ 個置く上記六基底で、$[\lambda_B\ \rho_B]$ の階数は18である。
ここで $\lambda_B,\rho_B$ は定理\ref{thm:ternary-ore-linear-space}の二写像である。
$\rho_B$ は階数15なので、初期補正の次元は高々15、六次元の左右乗法を割れば高々9となる。
\[
\begin{array}{c|c|c|c}
\lambda&\mu&\dim\ker(\lambda_B+\mu\rho_B)&
       \dim (P_B)_{\mu}\\ \hline
-1&1,-1&6,9&3,6\\
\sqrt{-1}&-1,\sqrt{-1},-\sqrt{-1}&6,6,6&3,3,3
\end{array}
\]
各固有核では $L_v-\mu R_v$ の三次元を除く。各核の残りが独立なので上界に達する。
さらに $c_j=0$ の頂点 $v=e_j^*$ ごとに、全固有核から
$d(v,v)\bmod kv$ への写像の階数は2である。コードは二次商の基底を
語空間へ持ち上げてこの二行を直接計算する。

特別な窓の $\mu$ 固有補正の代表を一つずつ取る。
離散付値環 $k[[s]]$ 上で、零だった $c_j$ は $s$、非零だったものはその値に取る。
位数2の場合は $b=1$ を保つ。位数4の場合は $b(1+b^2)/(1+b)=\prod_j c_j$
を $b(0)=-\lambda$ から解く。この点で左辺の導関数は非零なので解が存在する。
一般ファイバーは前節の同じ位数の滑らかな楕円基底である。
各次数の三生成代数は特殊・一般とも同じHilbert級数を持ち、各次数で有限生成の
離散付値環加群なので自由である。従って混合補正の固有核も自由加群の写像の核として
計算できる。前節と上表で一般・特殊の核次元が等しいので、単元小行列を分裂して
全ての特殊固有ベクトルを持ち上げられる。
その持上げを $\mu^n$ 倍して足せば全始点の混合三次条件が成立する。
定理\ref{thm:maximal-tangent-ore}の語の直和の証明はこの係数環上でも働くため、
四生成族は成分ごとに自由である。特殊補正が有限窓を与えれば、
射影的な各辺のファイバー次元の上半連続性から一般窓も四辺とも有限である。
反転で残る原始四乗根も従う。この平滑化は対称性を特別ファイバーに追加していない。
\end{proof}
